\begin{frame}{로지스틱회귀와 만나는 지점}
  \begin{block}{놀라운 사실 (Ng \& Jordan, 2001)}
    이렇게 구한 $P(y{=}1|x)$ 를 베이즈 정리로 전개하면, 정확히
    Week 2 로지스틱회귀의 \textbf{같은 시그모이드 형태}
    \[ P(y{=}1|x) = \sigma(w^\top x + b) \]
    가 나온다.
  \end{block}
  \vskip0.5em
  \begin{itemize}
    \item GDA = ``로지스틱회귀와 같은 결정 경계(직선)에 도달하는
      \textbf{또 다른 길}''.
    \item \textbf{트레이드오프}: 데이터가 실제로 가우시안에 가까우면
      GDA가 적은 데이터로도 잘 맞고, 가정이 틀렸으면 판별적
      (로지스틱회귀)이 더 안정적.
    \item ``모델을 데이터에 맞출 것인가, 데이터가 모델을 따르는가'' ---
      생성적 vs 판별적의 근본적 갈림.
    \item 아래 ``실현'' 절에서 수치로 확인: 2차원 가우시안 데이터에서
      GDA의 닫힌 형태 $w,b$ 와 경사하강 로지스틱 $w,b$ 는
      \textbf{거의 같은 방향}.
  \end{itemize}
\end{frame}
